
Autonomous agents and robots are increasingly expected to operate in complex real-world environments such as homes, warehouses, hospitals, and urban spaces. In these settings, agents must interact with diverse objects, perform varied tasks, and respond to changing environmental conditions. Unlike controlled laboratory settings, real-world environments are inherently open, new objects may appear, tasks may change, and interactions may produce outcomes that were not anticipated during system design or training. Developing systems that can operate robustly in such environments remains one of the central challenges in artificial intelligence and robotics.

Most contemporary AI and robotic systems are developed under closed-world assumptions, where the environment, tasks, and dynamics encountered during deployment closely match those observed during training. These assumptions simplify learning and planning but limit the ability of agents to operate reliably when confronted with novelty or unexpected situations. In open-world environments, previously learned models may become incomplete or invalid, action execution may fail due to unforeseen circumstances, and interaction policies may not generalize to new objects or tasks. Addressing these limitations requires systems that can adapt when prior knowledge becomes insufficient and acquire behaviors that generalize across diverse interactions with the environment.

This dissertation investigates how autonomous agents and robots can operate effectively in open-world environments through learning, adaptation, and generalization. The work focuses on three interconnected problems. First, developing environments and benchmarks that systematically introduce novelty and enable rigorous evaluation of open-world agents. Second, enabling agents to recover from failures during task execution by integrating planning with learning-based adaptation. Third, enabling robots to learn interaction policies that generalize across objects, tasks, and embodiment.

The first part of this dissertation addresses the need for systematic evaluation of open-world agents. We introduce benchmark environments that enable controlled injection of novelty and allow researchers to evaluate how agents respond to unexpected changes in their environment. These environments provide a platform for studying adaptive behavior and for comparing different approaches to open-world autonomy.

The second part of the dissertation focuses on adaptive planning-learning systems that allow agents to recover from execution failures. In open-world environments, plans may fail because of incomplete models or unexpected environmental changes. This work develops mechanisms that integrate task and motion planning with reinforcement learning-based recovery strategies, allowing agents to detect failures and adapt their behavior when previously learned knowledge becomes insufficient.

The third part of the dissertation investigates how robots can acquire generalizable interaction capabilities. Real-world robotic tasks often require sustained contact with objects and interactions that differ from those encountered during training. This work explores object-centric and force-based reinforcement learning approaches that enable robots to learn manipulation policies capable of generalizing across objects, tasks and particular robotic embodiments.

Taken together, the chapters of this dissertation develop a unified perspective on building autonomous agents and robots capable of operating in open-world environments. By combining benchmark-driven evaluation, adaptive planning-learning systems, and generalizable embodied interaction policies, this work advances the development of intelligent systems that can robustly operate in the real world.